
\def\PUDcodigo{ACE008}

\if\COMPILAR0
    \def\PUDcodacad{04507.8}
\else
    \def\PUDcodacad{04506.8}
\fi

\def\PUDdisciplina{Física 1}

\def\PUDsemestre{S2}

\def\PUDhoras{80}

%NOVOS ---------------------------------------------------
\def\PUDhorasTeoria{80}
\def\PUDhorasPratica{0}

\def\PUDinterECA{---}

\def\PUDatuaisECA{---}

\def\PUDrecursos{Projetor multimídia; Quadro branco e pincel.}

%----------------------------------------------------------

\def\PUDcreditos{4}

\def\PUDprerequisito{ \hyperref[PUD:ACE003]{ACE003} }

\if\COMPILAR0
    \def\PUDpreacad{ \hyperref[PUD:ACE003]{Cálculo 1 (04507.2)} }
\else
    \def\PUDpreacad{ \hyperref[PUD:ACE003]{Cálculo 1 (04506.2)} }
\fi

\def\PUDobjetivos{Compreender a física do movimento através dos seguintes objetivos específicos: Estruturar os conceitos das grandezas escalares e vetoriais e analisar os procedimentos das operações matemáticas entre vetores e suas aplicações na engenharia;  Formular conceitos físicos de deslocamento velocidade e aceleração através de seus respectivos gráficos interpretando os fenômenos mecânicos relacionados à cinemática dos corpos;  Identificar as componentes da velocidade e da aceleração de um corpo em movimento com trajetória parabólica e circular;  Analisar as leis de Newton dentro dos conceitos estáticos e dinâmicos aplicados em engenharia;  Compreender os efeitos ativo e passivo das forças de atrito em situações ligadas à engenharia;  Compreender o conceito de Trabalho de uma força, resolvendo problemas relacionados a potência e velocidade;  Comparar sistemas de forças conservativas e não conservativas, resolvendo problemas que envolvem energia mecânica em sistemas de forças
 gravitacionais e em sistemas de forças elásticas;  Calcular o centro de massa de um sistema de partículas e relacioná-lo com os problemas que envolvem Impulso e o Momento Linear.}

\def\PUDementa{Vetores.  Movimento retilíneo e no plano.  Leis de Newton.  Trabalho e energia.  Conservação da energia.  Centro de massa.  Rotação.  Momento linear.  Conservação do momento linear e colisões.}

\def\PUDconteudo{ &  Vetores: Grandezas escalares e vetoriais.  Vetor posição e deslocamento.  Representação geométrica das grandezas vetoriais.  Componentes vetoriais.  Método Analítico.  Operações com vetores (soma, subtração e multiplicação por um escalar).  Vetor unitário. 
\\ 
 &  Movimento em uma dimensão: Velocidade média.  Velocidade instantânea como derivada da posição.  Aceleração média.  Aceleração instantânea como derivada da velocidade.  Movimentos retilíneo uniforme e uniformemente variado.  Movimento vertical dos corpos. 
\\ 
 &  Movimento no Plano: Componentes ortogonais dos vetores: deslocamento, velocidade e aceleração.  Projéteis lançados horizontalmente: equações do movimento.  Projéteis lançados obliquamente: equações do movimento.  Movimento circular uniforme.  Posição, velocidade e aceleração relativas. 
\\ 
 &  Dinâmica da Partícula: Leis da Gravitação.  Primeira Lei de Newton.  Referenciais Inerciais.  Medida dinâmica da força.  Medida dinâmica da massa.  Segunda Lei de Newton.  Massa e peso.  Dinâmica no movimento circular uniforme.  Terceira Lei de Newton.  Medida estática da força.  Forças inerciais. 
\\ 
 &  Atrito: Coeficiente de atrito.  Forças de atrito. 
\\ 
 &  Trabalho e Energia: Operação com vetores.  Produto Escalar.  Trabalho de uma forca constante.  O Trabalho como a integral de uma força variável.  Teorema do Trabalho - Energia Cinética.  Potência. 
\\ 
 &  Conservação da Energia: Forças conservativas.  Forças não conservativas.  Energia Cinética, Energia Potencial e Energia Mecânica.  Lei da Conservação da Energia. 
\\ 
 &  Momento Linear e Colisões: Centro de massa.  Movimento do Centro de Massa.  Momento Linear.  Conservação do Momento Linear.  Impulso e Momento Linear.  Colisões. }

\def\PUDmetodo{Aulas expositórias que podem ser teóricas e/ou práticas, onde as práticas no laboratório serão marcadas no decorrer da disciplina.}

\def\PUDavalia{A avaliação será feita com aplicação de provas teóricas e/ou práticas, além da possibilidade de inclusão de trabalhos e seminários no decorrer da disciplina.}

\def\PUDbibbasica{ & \href{www.biblioteca.ifce.edu.br/index.asp?codigo_sophia=24032}{Resnick}, Robert.  Halliday, David, krane, Kenneth S.  Física 1.  Volume 1.  7º ed.  Rio de Janeiro, RJ : LTC, 2006.  356p.  ISBN: 9788521614845. 
\\ 
 & \href{www.biblioteca.ifce.edu.br/index.asp?codigo_sophia=65076}{Nussenzveig}, Herch M.  Curso de Física Básica -  Mecânica.  Volume 01.  5º ed.  São Paulo, SP: Edgard Blücher, 2013.  394p.  ISBN: 9788521207450. 
\\ 
 & \href{www.biblioteca.ifce.edu.br/index.asp?codigo_sophia=24335}{Sampaio}, José L.  Universo da Física I: Mecânica.  Volume 01.  2º ed.  Editora Atual.  São Paulo, 2005.  ISBN: 9788535705898. }

\def\PUDbibcomplementar{ & \href{www.biblioteca.ifce.edu.br/index.asp?codigo_sophia=24428}{Hugh} D.  Young e Roger A.  Freedman.  Física 1 - Mecânica.  12º ed.  Editora Pearson, 2008.  403p.  ISBN: 9788588639300. 
\\ 
 & \href{www.biblioteca.ifce.edu.br/index.asp?codigo_sophia=65137}{Alonso}, Marcelo;  Finn, Edward J.  Física 1: um curso universitário.  Volume 1.  2º ed.  São Paulo, SP : Edgard Blücher, 2014.  ISBN: 9788521208310. 
\\ 
 &   \href{www.biblioteca.ifce.edu.br/index.asp?codigo_sophia=61974}{Tipler}, Paul;  Mosca, Gene.  Fisica para cientistas e Engenheiros.  Volume 1.  6º ed.  Rio de Janeiro, RJ: LTC, 2016.  ISBN: 9788521617105.
\\
 &  \href{www.biblioteca.ifce.edu.br/index.asp?codigo_sophia=62872}{Hibbeler}, R. C.  Dinâmica: mecânica para engenharia.  12º ed.  São Paulo, SP: Pearson Prentice Hall, 2011.  591p.  ISBN: 9788576058144.
\\
 &  \href{www.biblioteca.ifce.edu.br/index.asp?codigo_sophia=8641}{Paraná}, Djalma Nunes.  Física. vol.1.  São Paulo, SP: Ática, 1998.  ISBN: 8508070802. }

\def\PUDautor{Samuel Vieira Dias}

\def\PUDdata{2013-04-23}

\def\PUDrevisao{1}

\def\PUDrevisor{Samuel Vieira Dias}

\def\PUDdatarevisao{2017-05-09}

\PUDcorpo
